\documentclass[11pt]{article}
\usepackage[margin=1in]{geometry}
\usepackage[T1]{fontenc}
\usepackage[utf8]{inputenc}
\usepackage{lmodern,microtype,amsmath,amssymb,amsthm,mathtools,booktabs}
\usepackage[colorlinks=true,linkcolor=blue,citecolor=blue,urlcolor=blue]{hyperref}
\hypersetup{pdftitle={The fixed-kernel regime of a margin-conditioned Bayesian neural-network prior},pdfauthor={Dmitry Vaintrob; Kaarel Hanni; OpenAI Codex}}
\setlength{\emergencystretch}{2em}
\newtheorem{theorem}{Theorem}[section]
\newtheorem{proposition}[theorem]{Proposition}
\newtheorem{lemma}[theorem]{Lemma}
\newtheorem{corollary}[theorem]{Corollary}
\theoremstyle{remark}
\newtheorem{remark}[theorem]{Remark}
\DeclareMathOperator*{\argmin}{arg\,min}
\title{The fixed-kernel regime of a margin-conditioned\\Bayesian neural-network prior}
\author{Dmitry Vaintrob \and Kaarel H\"anni \and OpenAI Codex}
\date{October 15, 2026}
\begin{document}
\maketitle
\begin{abstract}
We study the width limit of a fully connected $\tanh$ network on the
Boolean cube with an appended constant input, under the Gaussian
scaling $\operatorname{Var}(W^1_{ij})=N^{1-\gamma}/d$,
$\operatorname{Var}(W^s_{ij})=N^{-\gamma}$ for later hidden layers, and
$\operatorname{Var}(W^L_i)=N^{-\gamma-1}$ for the readout. The Bayesian
prior is conditioned on an output margin at every Boolean input;
training then conditions on observed signs. For every fixed
$\gamma<1$, dimension, and depth, hidden-feature covariances concentrate
at speed $N$ around an explicit iterated-arcsine kernel. For
$0<\gamma<1$, outputs satisfy a Gaussian large-deviation principle at
speed $N^\gamma$, and the actual posterior concentrates on compatible
Boolean functions of minimum kernel margin energy. At $\gamma=0$ the
limit is a Gaussian-process classifier with the global margin
restriction; at $\gamma<0$ the margin disappears and the limit is a
Gaussian sign classifier. These establish a fixed-kernel regime under
the original prediction conventions, while distinguishing it from
Gaussian-process regression. We derive exact parity prior rates, a
strictly paired kernel spectrum, the complete set of unconditioned
minimum-energy sign functions, and classification lower bounds. A
prediction counts as successful only when the posterior probability
of its correct label exceeds $2/3$. For $\gamma\le0$ and depth
at least a constant multiple of $d$, we prove that, for every target, the
failure fraction is exactly the unseen fraction in the width limit; even
parity on $d$ bits then needs $\Theta(2^d)$ examples. The
stronger claim that every fixed
$\gamma<1$ requires exponentially many examples for a $k$-bit parity
is not established: a small parity eigenvalue alone does not imply
that claim for this posterior.
Note added (October 2026): the later paper \emph{Exponential parity sample complexity
below the mean-field scaling} proves an $\exp(\Omega(k))$ confidence lower bound
for full $k$-bit parity ($d=k$) at every fixed $\gamma<1$ and fixed depth, in the
width-first limit. For sparse parity with $0<\gamma<1$, the paper \emph{Exponential sample
lower bounds for sparse parity under a globally separated neural prior} proves an
$\exp(\Omega(k))$ bound; the natural $\exp\{\Omega(k\log(ed/k))\}$ scale at fixed depth
remains open.
\end{abstract}

\section{Model, posterior, and order of limits}

Fix $d\ge1$, depth $L\ge2$, margin $a>0$, and a real exponent
$\gamma<1$. Inputs are $x\in\mathcal X=\{-1,1\}^d$, with
$h^0(x)=(1,x)$. All hidden layers have width $N$ and
\[
 h^s(x)=\tanh(W^s h^{s-1}(x)),\quad 1\le s\le L-1,
 \qquad z_N(x)=W^Lh^{L-1}(x).
\]
The activation acts coordinatewise. There are no biases after the
appended constant input. All weights are independent and centered
Gaussian, with variances
\begin{equation}\label{eq:model-variances}
 \operatorname{Var}(W^1_{ij})=\frac{N^{1-\gamma}}d,
 \qquad \operatorname{Var}(W^s_{ij})=N^{-\gamma}\ (1<s<L),
 \qquad \operatorname{Var}(W^L_i)=N^{-\gamma-1}.
\end{equation}
This is the strong-prior scaling line in the source notes~\cite{notes},
continued below $\gamma=1$. The main statements are specific to this
line. Changing the relative layer variances or inserting additional
biases can change the limiting kernel.

Let $Z_N=(z_N(x))_{x\in\mathcal X}$. The prior on Boolean functions is
obtained by first conditioning on
\[
 A_a=\{|Z_N(x)|>a\text{ for all }x\in\mathcal X\}
\]
and then taking $g(x)=\operatorname{sign}Z_N(x)$. A training sample
$S=((X_i,f(X_i)))_{i=1}^n$ supplies noiseless signs. Write
$\mathcal G_S=\{g:g(X_i)=f(X_i)\text{ for all }i\}$. For
$A_g=\{g(x)Z_N(x)>a\text{ for every }x\}$, the exact posterior is
\begin{equation}\label{eq:exact-finite-posterior}
 \pi_N(g\mid S)=
 \frac{\Pr(A_g)}{\sum_{h\in\mathcal G_S}\Pr(A_h)}
 \quad(g\in\mathcal G_S),
\end{equation}
with zero mass outside $\mathcal G_S$. The normalization introduced by
the initial global-margin conditioning cancels in this formula.

At an input $x$, let
$p_N(x\mid S,f)=\pi_N\{g:g(x)=f(x)\mid S\}$. We count a prediction as
successful only when $p_N(x\mid S,f)>2/3$. Our primary error criterion
is therefore the fraction of inputs failing this confidence requirement:
\begin{equation}\label{eq:confidence-error}
 F_N(S,f)=\frac1{2^d}\sum_{x\in\mathcal X}
                    \mathbf 1\{p_N(x\mid S,f)\le2/3\}.
\end{equation}
As a useful auxiliary quantity, define the average posterior
probability of a wrong label:
\begin{equation}\label{eq:posterior-error}
 R_N(S,f)=\frac1{2^d}\sum_{x\in\mathcal X}
            \pi_N\{g:g(x)\ne f(x)\mid S\}.
\end{equation}
This is also the expected test error of a Boolean function drawn from
the posterior. Pointwise comparison of the wrong-label probability
on the failing and successful inputs gives
\begin{equation}\label{eq:confidence-risk-comparison}
 \max\{0,(3R_N-1)/2\}\le F_N\le\min\{1,3R_N\}.
\end{equation}
Indeed, failing inputs each contribute at least $1/3$ to $R_N$,
while successful inputs contribute less than $1/3$ and all inputs
contribute at most one. The same linear bounds hold after averaging
over samples. In particular, posterior error near $1/2$ forces
confidence failure on at least about a quarter of inputs.

The majority rule chooses the more probable sign at
each input and is a different performance criterion. Its error is at
most twice~\eqref{eq:posterior-error}; the reverse implication fails.
A posterior that assigns probability $0.5001$ to every correct sign
can have a perfect majority classifier and posterior error close to
$1/2$.

Every width limit in this paper takes $N\to\infty$ with $d,L,\gamma,a$
fixed. Sample-complexity statements then study the resulting family
as $d,k,n$ vary, where $k$ is the parity size. We do not exchange these
limits or claim quantitative uniformity in $d$ or $L$.

\section{The regime \texorpdfstring{$\gamma<1$}{gamma < 1}: a frozen kernel and three posterior limits}
\label{sec:gamma-less-one}

Throughout this section the activation is $\tanh$, the depth $L\ge2$ and
input dimension $d\ge1$ are fixed, and the width tends to infinity. Write
$\mathcal X=\{-1,1\}^d$, $M=2^d$, and let $a>0$ be the required margin.
All limits below hold at each fixed $\gamma<1$; uniformity in dimension,
depth, or $\gamma\uparrow1$ is not asserted. Define
\begin{equation}\label{eq:arcsine-kernel}
 K^{(0)}(x,y)=\frac{1+x\cdot y}{d+1},\qquad
 K^{(s+1)}(x,y)=\frac2\pi\arcsin K^{(s)}(x,y),\qquad
 K=K^{(L-1)}.
\end{equation}
Thus the limiting kernel uses the saturated activation. It is not the
unit-variance, finite-gain $\tanh$ kernel.

\begin{theorem}[Exponential freezing of the feature kernel]
\label{thm:kernel-freezing}
Let $h_i^s(x)$ denote the activation of hidden neuron $i$ in layer $s$,
and set
\[
 Q_N^s(x,y)=\frac1N\sum_{i=1}^N h_i^s(x)h_i^s(y).
\]
For every $\epsilon>0$ and each $1\le s\le L-1$, there are constants
$c,C>0$ and $N_0$ such that
\[
 \Pr\!\left\{\max_{x,y}|Q_N^s(x,y)-K^{(s)}(x,y)|>\epsilon\right\}
 \le C e^{-cN}\qquad(N\ge N_0).
\]
The matrix $K$ is positive definite. Conditional on all hidden layers,
the output vector $Z_N=(z_N(x))_{x\in\mathcal X}$ has the exact law
\begin{equation}\label{eq:conditional-output-kernel}
 Z_N\mid (\text{hidden layers})\sim
 \mathcal N(0,N^{-\gamma}Q_N^{L-1}).
\end{equation}
Consequently $N^{\gamma/2}Z_N$ converges in distribution to
$G\sim\mathcal N(0,K)$.
\end{theorem}

\begin{proof}
Put $b_N=N^{(1-\gamma)/2}\to\infty$. Conditional on the preceding
layer, different preactivation rows are independent Gaussian vectors
with covariance $b_N^2Q_N^{s-1}$. For the first layer, replace this
covariance matrix by the deterministic matrix
$b_N^2R_0$, where $R_0(x,y)=(1+x\cdot y)/d$; its correlation matrix is
$K^{(0)}$.

If $U\sim\mathcal N(0,q)$, then bounding the Gaussian density by its
maximum and integrating $1-\tanh v$ gives
\[
 \mathbb E|\tanh(bU)-\operatorname{sign}U|
 \le \frac{2\log2}{b\sqrt{2\pi q}}.
\]
For a centered Gaussian pair with correlation $r$, the expected product
of its signs is $2\arcsin(r)/\pi$. Thus, uniformly when coordinate
variances are bounded below, the expected product of the two saturated
activations differs from this arcsine expression by $O(b_N^{-1})$.
The appended constant input implies
$|K^{(s)}(x,y)|<1$ for distinct $x,y$. In a sufficiently small fixed
neighborhood of $K^{(s)}$, the map sending a covariance matrix to its
normalized arcsine matrix is Lipschitz, and all diagonal variances are
bounded below. Finally, conditional Hoeffding bounds for the $N$
independent products in $[-1,1]$ give
\[
 \Pr\!\left\{\max_{x,y}|Q_N^s(x,y)
 -\mathbb E[Q_N^s(x,y)\mid\text{preceding layer}]|>t
 \,\middle|\,\text{preceding layer}\right\}
 \le 2M^2 e^{-Nt^2/2}.
\]
Induction over the fixed number of layers proves the concentration
statement, absorbing the deterministic $O(b_N^{-1})$ error by taking
$N$ sufficiently large.

For positive definiteness, $K^{(1)}$ is the Gram matrix of the functions
$u\mapsto\operatorname{sign}(u\cdot(1,x))$ under a nondegenerate
Gaussian measure. They are linearly independent: in a vanishing linear
combination, cross the hyperplane $u\cdot(1,x)=0$ at a point on none of
the other finitely many hyperplanes. Only the corresponding term
changes sign, so its coefficient is zero. Such a point exists because
the augmented vectors $(1,x)$ are pairwise nonparallel. The positive
power-series coefficients of the arcsine function and the Schur product
theorem then imply $K^{(s+1)}\succeq(2/\pi)K^{(s)}$, preserving strict
positive definiteness. The independent Gaussian readout proves
\eqref{eq:conditional-output-kernel}, and convergence follows.
\end{proof}

\subsection{Positive \texorpdfstring{$\gamma$}{gamma}: the actual posterior is a kernel-energy learner}

For a Boolean function $g$ define
\begin{equation}\label{eq:kernel-orthant-energy}
 E_K(g)=\min\left\{\frac12z^TK^{-1}z:
               g(x)z(x)\ge a\text{ for every }x\in\mathcal X\right\}.
\end{equation}
This minimum is finite and has a unique minimizing vector, by strict
convexity and coercivity.

\begin{theorem}[Kernel posterior limit]
\label{thm:kernel-posterior}
For each fixed $0<\gamma<1$, $Z_N$ satisfies a large-deviation principle
with speed $N^\gamma$ and good rate function
$I(z)=z^TK^{-1}z/2$. In particular,
\begin{equation}\label{eq:kernel-orthant-rate}
 \lim_{N\to\infty}-N^{-\gamma}
 \log\Pr\{g(x)Z_N(x)>a\text{ for every }x\}=E_K(g).
\end{equation}
Let $\mathcal G_S$ be the Boolean functions consistent with the observed
training labels and let $E_S=\min_{h\in\mathcal G_S}E_K(h)$. Under the
all-input-margin prior followed by noiseless sign conditioning,
\begin{equation}\label{eq:kernel-posterior-rate}
 \lim_{N\to\infty}-N^{-\gamma}
 \log\Pr\{\operatorname{sign}Z_N=g\mid\text{all margins and }S\}
 =E_K(g)-E_S\qquad(g\in\mathcal G_S).
\end{equation}
The posterior therefore concentrates on the functions in
$\operatorname*{argmin}_{g\in\mathcal G_S}E_K(g)$.
\end{theorem}

\begin{proof}
On a common probability space let
\[
 Z_N=N^{-\gamma/2}(Q_N^{L-1})^{1/2}\xi,
 \qquad Y_N=N^{-\gamma/2}K^{1/2}\xi,
 \qquad \xi\sim\mathcal N(0,I_M),
\]
with $\xi$ independent of the hidden layers. The square-root map is
locally Lipschitz near the positive definite $K$. For fixed
$\delta,\epsilon>0$, Theorem~\ref{thm:kernel-freezing} consequently gives
\[
 \Pr\{\|Z_N-Y_N\|>\delta\}
 \le C_\epsilon e^{-c_\epsilon N}
   +\Pr\!\left\{\|\xi\|>
          \frac{\delta N^{\gamma/2}}{C\epsilon}\right\}.
\]
Divide logarithms by $N^\gamma$, first let $N\to\infty$, and then let
$\epsilon\downarrow0$. Because $\gamma<1$, the resulting upper limit
is $-\infty$. Thus $Z_N$ is exponentially equivalent, at speed
$N^\gamma$, to the elementary small-noise Gaussian family $Y_N$.
The latter has rate $I$, proving the large-deviation statement.

The strict margin orthant is open and its closure is the feasible set
in~\eqref{eq:kernel-orthant-energy}. Scaling its unique minimizer by
$1+\eta$ puts it in the open orthant and approaches the same energy
as $\eta\downarrow0$. The open lower bound and closed upper bound
therefore match. Summing over the finitely many compatible sign
patterns and subtracting logarithmic asymptotics proves
\eqref{eq:kernel-posterior-rate}.
\end{proof}

If the minimizing Boolean function is unique, the limiting posterior
is its point mass. If there is a tie, the exponential rates do not
determine the mixture weights. The result does not prescribe a uniform
mixture of minimizers.

The limit is governed by a fixed kernel, but its optimization problem
retains the constraint $|z(x)|\ge a$ at every unseen input. Its feasible
set is
\[
 \{z:f(x_i)z(x_i)\ge a\text{ at observed inputs},\quad
             |z(x)|\ge a\text{ at every input}\}.
\]
This is a nonconvex union of orthants. Dropping the unseen-input
constraints produces a different learner. In particular, ordinary
kernel interpolation, Gaussian-process regression, and the usual
hard-margin support-vector machine do not describe this posterior.

\subsection{Zero and negative \texorpdfstring{$\gamma$}{gamma}}

\begin{proposition}[The remaining fixed-kernel limits]
\label{prop:nonpositive-gamma}
At $\gamma=0$, the posterior on Boolean functions converges to that of
$G\sim\mathcal N(0,K)$ conditioned on $|G(x)|>a$ at every input and on
the observed training signs. At $\gamma<0$, it converges to the sign
law of the same Gaussian vector conditioned only on observed training
signs.
\end{proposition}

\begin{proof}
The normalized outputs converge weakly to the nondegenerate $G$.
At $\gamma=0$, every relevant conditioning event has boundary in a
finite union of hyperplanes of Gaussian probability zero, and its
probability is strictly positive. Conditional probabilities converge.
At $\gamma<0$, the normalized margin is
$aN^{\gamma/2}\to0$. The probability that any normalized coordinate
falls within this margin tends to zero. Every prescribed finite set
of observed signs has positive limiting probability, so the same
conditional convergence argument applies.
\end{proof}

\subsection{What the kernel limit says about parity}

Cube symmetries diagonalize $K$ in the parity basis. If
$\chi_U(x)=\prod_{i\in U}x_i$ and $\lambda_U$ is its eigenvalue for
the normalized operator $M^{-1}K$, then
\begin{equation}\label{eq:parity-kernel-energy}
 E_K(\chi_U)=\frac{a^2}{2\lambda_U}.
\end{equation}
Indeed $a\chi_U$ is feasible, and its gradient
$K^{-1}(a\chi_U)$ is a positive multiple of $\chi_U$. Its inner product
with every feasible displacement is nonnegative; the first-order
criterion for convex minimization proves optimality.

Equation~\eqref{eq:parity-kernel-energy} is an exact prior-rate
calculation. It is not a sample-complexity lower bound. For
$0<\gamma<1$, that question requires comparison with the energies of
all compatible Boolean functions; for $\gamma=0$ it concerns globally
margin-conditioned Gaussian orthant probabilities; for $\gamma<0$
it concerns zero-margin Gaussian sign probabilities. A small parity
eigenvalue by itself does not answer any of these posterior-learning
questions. The regression subspace argument also does not apply: the
full-cube conditioning need not leave the posterior mean in the span
of training kernel sections. Moreover a lower bound stated in the
parity size $k$ must specify the ambient dimension, depth, input
distribution, order of limits, and whether error means posterior
misclassification probability or majority-vote error.


\section{Classification bounds and the limits of eigenvalue arguments}
\label{sec:classification-bounds}

\subsection{The spectrum and the first functions eliminated by data}

\begin{proposition}[Paired spectrum and minimum-energy signs]
\label{prop:paired-spectrum}
Let $\lambda_r$ be the eigenvalue of $M^{-1}K$ on degree-$r$ parities.
For the kernel~\eqref{eq:arcsine-kernel},
\[
 \lambda_0=\lambda_1>\lambda_2=\lambda_3>
 \lambda_4=\lambda_5>\cdots>0,
\]
with a final unpaired $\lambda_d$ when $d$ is even. The minimum of
$E_K(g)$ over all Boolean functions is $a^2/(2\lambda_0)$, attained
exactly by the constants and signed coordinate functions
$g\in\{1,-1,x_1,-x_1,\ldots,x_d,-x_d\}$.
For every finite width, all these functions have identical
unconditioned full-margin prior probabilities.
\end{proposition}

\begin{proof}
Put $D=d+1$ and write $F$ for the $(L-1)$-fold composition of
$t\mapsto2\arcsin(t)/\pi$. Its power series is
$F(t)=\sum_{m\text{ odd}}c_mt^m$, with every $c_m>0$ and
$\sum_m c_m=1$. Positivity follows from the arcsine series and is
preserved under composition. First consider the kernel $F(u\cdot v/D)$
on the full $D$-dimensional cube. The degree-$r$ Fourier eigenvalue
of $(u\cdot v/D)^m$ is the probability $p_m(r)$ that a simple random
walk, starting at one cube vertex and flipping a uniformly chosen
coordinate at each step, reaches a specified vertex at Hamming
distance $r$ after $m$ steps. This follows by expanding the power:
a term survives Fourier averaging exactly when the coordinates
chosen an odd number of times form the specified support.

These probabilities decrease with distance within each parity class.
Indeed, for $0\le r\le D-2$,
\begin{align*}
 D\bigl[p_{m+1}(r)-p_{m+1}(r+2)\bigr]
 ={}&r\bigl[p_m(r-1)-p_m(r+1)\bigr]\\
 &+(D-r-2)\bigl[p_m(r+1)-p_m(r+3)\bigr],
\end{align*}
where terms with zero coefficients at the boundaries are omitted.
Induction starts from the point mass at time zero. Since
$p_r(r)>0=p_r(r+2)$ and $c_r>0$ for odd $r$, the positive odd-degree
eigenvalues of the full-cube kernel decrease strictly; its even-degree
eigenvalues vanish. Restricting the distinguished coordinate to $+1$
makes the reduced degree-$r$ eigenvalue the sum of the full-cube
degree-$r$ and degree-$(r+1)$ eigenvalues. This gives the paired
ordering stated.

The largest matrix eigenvalue of $K$ is $M\lambda_0$. Every
full-margin output satisfies
\[
 \frac12z^TK^{-1}z\ge\frac{\|z\|_2^2}{2M\lambda_0}
                     \ge\frac{a^2}{2\lambda_0}.
\]
Equality requires both $|z(x)|=a$ everywhere and
$z/a$ in the span of $1,x_1,\ldots,x_d$. An affine function
$c+\sum_j b_jx_j$ is Boolean only if it is a constant or signed
coordinate: expanding its square shows $b_ib_j=0$ for $i\ne j$
and $cb_i=0$, so at most one coefficient is nonzero, with magnitude
one. These functions attain equality, completing the characterization.

To prove the finite-width equality of prior weights, exchange
columns $0$ and $j$ of the first-layer weight matrix. Its Gaussian
law is unchanged. Let $P_j$ denote this exchange and let $T_j$ be
the cube bijection defined by $(T_jx)_j=x_j$ and
$(T_jx)_i=x_jx_i$ for $i\ne j$. Then
$P_j(1,x)=x_j(1,T_jx)$. The network as a function of its augmented
input is odd, since all activations are odd and there are no later
biases. Thus the transformed output is $x_j z_N(T_jx)$.
The all-positive margin orthant is mapped bijectively to that of
$x_j$, proving equality. Flipping the readout weights proves
equality with the negative functions as well.
\end{proof}

\begin{corollary}[A posterior-error lower bound from surviving coordinates]
\label{cor:surviving-coordinates}
Fix $0<\gamma<1$ and let $f=\chi_U$ be a parity of size $k\ge2$.
If any constant or signed coordinate agrees with a training sample
$S$, the posterior converges to the uniform distribution over all
such compatible functions. In particular,
\[
 \lim_{N\to\infty}R_N(S,f)=\frac12.
\]
For $d>k$ and $n\ge1$ independent uniform training examples,
\begin{equation}\label{eq:coordinate-sample-bound}
 \liminf_{N\to\infty}\mathbb E_S R_N(S,f)
 \ge\frac12\left[1-\left(1-2^{1-n}\right)^{d-k}\right].
\end{equation}
\end{corollary}

\begin{proof}
If such a function survives, the compatible absolute minimizers in
Proposition~\ref{prop:paired-spectrum} exhaust the compatible
minimum-energy functions. Theorem~\ref{thm:kernel-posterior} forces
the posterior onto these functions, and their exact equal prior
weights make the conditional distribution uniform. Each has uniform error exactly
$1/2$ against a parity of degree at least two, independently of
the mixture weights. This proves the first statement.

Conditional on all samples in coordinates belonging to $U$, the
target labels are fixed. Each of the $d-k$ irrelevant coordinates
has an independent uniform sample column in $\{-1,1\}^n$. One of
its two signs agrees with all labels with probability $2^{1-n}$.
These events are conditionally independent. The probability that
at least one such coordinate survives is consequently
$1-(1-2^{1-n})^{d-k}$. On this event the limiting error is $1/2$;
on its complement error is nonnegative. Fatou's lemma proves the
bound.
\end{proof}

\begin{lemma}[Parity correlation with a halfspace]
\label{lem:halfspace-parity}
If $h$ is a Boolean halfspace on $\{-1,1\}^d$ and $|U|=k\ge1$, then
$|\mathbb E[\chi_U h]|\le1/\sqrt{k}$. Either inclusion or exclusion
of the threshold boundary is allowed.
\end{lemma}

\begin{proof}
Condition on coordinates outside $U$ and flip selected input
coordinates so that the halfspace is monotone in the remaining $k$
coordinates. These flips change the parity only by a constant sign.
For a monotone Boolean function $h$, its discrete derivative
$D_i h=(h(x_i=1)-h(x_i=-1))/2$ is nonnegative. Thus, for every
$i\in U$,
\[
 |\widehat h(U)|
 =|\mathbb E[\chi_{U\setminus\{i\}}D_ih]|
 \le\mathbb E D_i h=\widehat h(\{i\}).
\]
Summing over the $k$ coordinates and using Cauchy--Schwarz and
Parseval gives $k|\widehat h(U)|\le\sqrt{k}$.
Average the conditional estimate over the other coordinates.
\end{proof}

\begin{corollary}[Failure to attain posterior confidence $2/3$]
\label{cor:confidence-groundstates}
Under the hypotheses of Corollary~\ref{cor:surviving-coordinates},
if any constant or signed coordinate is consistent with $S$, then
\[
 \liminf_{N\to\infty}F_N(S,f)
 \ge b_k:=\max\left\{\frac14,\frac12-\frac1{2\sqrt{k}}\right\}.
\]
If there are $m$ such compatible functions, the additional bound
$\liminf_N F_N(S,f)\ge1-9/m$ holds.
Consequently, when $d>k$ and $n\ge1$,
\[
 \liminf_{N\to\infty}\mathbb E_S F_N(S,f)
 \ge b_k\left[1-\left(1-2^{1-n}\right)^{d-k}\right].
\]
\end{corollary}

\begin{proof}
The bound $1/4$ follows from
\eqref{eq:confidence-risk-comparison} and
Corollary~\ref{cor:surviving-coordinates}. To obtain the stronger
bound for large $k$, write $\epsilon_N\to0$ for the posterior mass
outside the consistent absolute minimizers. The mean Boolean
prediction conditioned on those minimizers is an affine function
$\widetilde m_N(x)=c_N+\sum_jb_{N,j}x_j$. The full posterior mean
$m_N(x)$ satisfies
$\|m_N-\widetilde m_N\|_\infty\le2\epsilon_N$.
Confidence success is the event $f(x)m_N(x)>1/3$, which implies
$f(x)\widetilde m_N(x)>t_N:=1/3-2\epsilon_N>0$ for large $N$.

For any affine $m$ and any $t>0$, set
$w(x)=\mathbf1\{m(x)>t\}-\mathbf1\{m(x)<-t\}$.
This is the average of two Boolean halfspaces, using the appropriate
boundary conventions. By Lemma~\ref{lem:halfspace-parity},
$|\mathbb E fw|\le1/\sqrt{k}$. Moreover
\[
 \Pr\{fm>t\}=\frac12\mathbb E(|w|+fw)
          \le\frac12+\frac1{2\sqrt{k}}.
\]
Apply this to $m=\widetilde m_N$ and $t=t_N$. This proves the
pointwise-in-sample bound without assuming convergence of the
weights among tied minimizers. The expected bound follows using
the event probability calculated in
Corollary~\ref{cor:surviving-coordinates} and Fatou's lemma.

For the additional bound, equal weights make
$\widetilde m_N=m^{-1}\sum_{j=1}^m g_j$ independent of $N$.
Distinct signed constants and coordinates are orthogonal unless
they are negatives of one another. Hence
$\mathbb E\widetilde m_N^2\le1/m$. Confidence success implies
$|\widetilde m_N|>t_N$, so Markov's inequality bounds its fraction
by $1/(mt_N^2)\to9/m$.
\end{proof}

More explicitly, put $D=d-k$ and $p=2^{1-n}$ for $n\ge1$.
The number of surviving irrelevant coordinates is
$\operatorname{Bin}(D,p)$. For any integer $t>9$, the same proof gives
\[
 \liminf_N\mathbb E_S F_N(S,f)
 \ge (1-9/t)\Pr\{\operatorname{Bin}(D,p)\ge t\}.
\]
Thus confidence failure can be close to one while many spurious
coordinates remain compatible with the sample.

These results establish a necessary sample scale of order
$\log(d-k)$ for posterior confidence in this regime. They are not
exponential-in-$k$ bounds, nor do they assert the same error for
posterior majority.
The spectrum also gives
$E_K(\chi_U)\ge(a^2/2)\binom dk$, since
$\sum_{r=0}^d\binom dr\lambda_r=1$. Even an exponentially large
parity energy when $d\ge2k$ is a prior statement rather than a
sample lower bound.


\subsection{An information bound that does apply to the posterior}

\begin{proposition}[Parity-identification bound]
\label{prop:information-bound}
Let $1\le k\le d$ and $B=\binom dk$. For $n$ independent uniform
training inputs and any fixed degree-$k$ parity $f$, the posterior
learner in this paper satisfies
\begin{equation}\label{eq:information-lower-bound}
 \mathbb E_S R_N(S,f)\ge
 \frac12\left(1-\sqrt{\frac{2^n}{B}}\right)
\end{equation}
whenever its posterior is defined. The same statement holds for its
width limits and for a majority classifier whose tie rule respects
coordinate permutations. Thus expected error at most
$\epsilon<1/2$ requires
\[
 n\ge\log_2\binom dk+2\log_2(1-2\epsilon).
\]
\end{proposition}

\begin{proof}
First average the target uniformly over the $B$ degree-$k$ parities
and fix the training inputs. The targets split into at most $2^n$
cells $C_y$ according to their label vectors. Let $m_y(x)\in[-1,1]$
denote the learner's mean signed prediction given the label vector
$y$. For the posterior learner it is the posterior marginal mean of
$g(x)$; for majority prediction it is the selected sign. Use the
uniform inner product on the Boolean cube. Since
$\|m_y\|_2\le1$, orthogonality of the parities and Bessel's inequality
give
\[
 \sum_{U\in C_y}\langle m_y,\chi_U\rangle
 \le \sqrt{|C_y|}
       \left(\sum_{U\in C_y}\langle m_y,\chi_U\rangle^2\right)^{1/2}
 \le \sqrt{|C_y|}.
\]
Summing over cells and applying Cauchy--Schwarz bounds average target
correlation by $\sqrt{2^n/B}$. Classification error is half of one
minus this correlation. Finally, coordinate permutations act
transitively on degree-$k$ supports and preserve the network prior,
the all-input margin event, the sampling distribution, and the
posterior rule. The expected error is therefore identical for every
degree-$k$ parity. This proves the fixed-target result as well as its
stated extensions.
\end{proof}

This lower bound is of order $k\log(d/k)$ when $k\ll d$, and is
vacuous when $k=d$. It is not exponential in $k$. It also illustrates
why the ambient dimension matters in any statement about ``an XOR of
size $k$.''
For the primary confidence criterion, the same argument and
\eqref{eq:confidence-risk-comparison} imply
\[
 \mathbb E_S F_N(S,f)\ge
 \max\left\{0,\frac14-\frac34\sqrt{\frac{2^n}{\binom dk}}\right\}.
\]

\subsection{Exponential confidence sample complexity at sufficient depth}

We can prove exponential sample complexity for the actual scaling
family when $\gamma\le0$ and depth is allowed to grow linearly with
dimension. This uses a uniform comparison after conditioning on
all margins and any training signs. Its uniformity is essential:
the conditioning event can itself be exponentially unlikely.

\begin{lemma}[Gaussian comparison after margins and signs]
\label{lem:gaussian-comparison}
For each $a\ge0$ there is $c_a>0$ such that, for an $M\times M$
covariance matrix $K$ and $0<\eta<1/4$, the condition
\[
 \|K-I\|_{\mathrm{op}}\le c_a\eta/M
\]
implies the following. If $B$ prescribes $|z_j|>a$ at every coordinate
and arbitrary signs on any subset, then
\[
 \operatorname{TV}\bigl(\mathcal N(0,K)(\cdot\mid B),
                        \mathcal N(0,I)(\cdot\mid B)\bigr)\le\eta.
\]
When $a=0$, the all-coordinate absolute-value condition is vacuous.
\end{lemma}

\begin{proof}
Let $Q=\mathcal N(0,I)(\cdot\mid B)$,
$H=K^{-1}-I$, and $h=\|H\|_{\mathrm{op}}$. Conditional densities have
ratio $w/\mathbb E_Qw$, where $w(z)=\exp(-z^THz/2)$.
Under $Q$, the squared coordinates are independent with the law of
$G^2$ conditional on $|G|>a$, independently of the prescribed signs.
Put $\mu_a=\mathbb E[G^2\mid |G|>a]$. For a finite $D_a$,
\[
 \log\mathbb E[e^{tG^2}\mid |G|>a]\le D_at
 \quad(0\le t\le1/4),
\]
because the moment-generating function and its derivative are finite
on this interval. With $T=\sum_jz_j^2$ and $h\le1/2$,
\[
 \mathbb E_Q|w-1|\le e^{D_aMh/2}-1,
 \qquad \mathbb E_Qw\ge e^{-\mu_aMh/2}.
\]
The second inequality uses Jensen's inequality. Thus
\[
 \operatorname{TV}
 \le\frac{\mathbb E_Q|w-1|}{\mathbb E_Qw}
 \le e^{\mu_aMh/2}(e^{D_aMh/2}-1).
\]
For $Mh\le1$ the last expression is at most a constant depending
only on $a$ times $Mh$. The inverse identity gives
$h\le2\|K-I\|_{\mathrm{op}}$ when the latter norm is at most $1/2$.
Choose $c_a$ small enough to obtain the claim.
\end{proof}

\begin{lemma}[Decay of saturated-kernel correlations]
\label{lem:arcsine-decay}
There is an absolute constant $C$ such that the kernel in
\eqref{eq:arcsine-kernel} obeys $K(x,x)=1$ and
$|K(x,y)|\le\rho$ for every distinct $x,y$, whenever
\[
 L-1\ge C\bigl(\log(d+1)+\log(1/\rho)+1\bigr).
\]
\end{lemma}

\begin{proof}
Write $T(r)=2\arcsin(r)/\pi$. This map is odd and satisfies
$0<T(r)<r$ for $0<r<1$. Since $T'(r)\to\infty$ as $r\uparrow1$,
there is $r_*<1$ for which
$1-T(r)\ge2(1-r)$ whenever $r_*\le r<1$.
Initially the absolute correlation of any distinct augmented cube
points is at most $(d-1)/(d+1)$. Its distance from one therefore
doubles until its iterates fall below $r_*$, requiring at most
$O(\log(d+1)+1)$ steps. The positive power series shows
$T(r)/r\le T(r_*)/r_*=:q<1$ for $0<r\le r_*$.
A further $O(\log(1/\rho)+1)$ steps suffice. Oddness handles
negative correlations, and $T(1)=1$ preserves the diagonal.
\end{proof}

\begin{theorem}[Exact confidence failure in a deep subfamily]
\label{thm:deep-confidence}
Fix $\gamma\le0$ and $a>0$. There is a constant $C_a$ such that,
if $L\ge C_a(d+1)$, then for every target and every fixed training
sample $S$, the width-limit confidence failure is exactly
\[
 \lim_{N\to\infty}F_N(S,f)=\frac{U_S}{2^d},
\]
where $U_S$ is the number of cube inputs absent from the sample.
Consequently, for $n$ independent uniform training inputs,
\begin{equation}\label{eq:deep-exact-confidence}
 \lim_{N\to\infty}\mathbb E_S F_N(S,f)
       =(1-2^{-d})^n.
\end{equation}
For $0<\epsilon\le1/2$, expected confidence failure at most
$\epsilon$ therefore requires and is achieved by
$\Theta(2^d\log(1/\epsilon))$ examples. In particular, for $d=k$
this is exponential in the size of a parity circuit.
\end{theorem}

\begin{proof}
Choose a fixed $\eta<1/6$, for example $\eta=1/12$.
Apply Lemma~\ref{lem:arcsine-decay} with
$\rho=c\eta/M^2$, where $M=2^d$ and $c$ is the smaller of the
comparison constants for margins $a$ and zero. Its depth requirement
is $O(d+1)$, and
$\|K-I\|_{\mathrm{op}}\le(M-1)\rho\le c\eta/M$.
At $\gamma=0$ use Lemma~\ref{lem:gaussian-comparison} with margin
$a$; at $\gamma<0$ use it with margin zero, in accordance with
Proposition~\ref{prop:nonpositive-gamma}.

Under the independent Gaussian comparator, every unobserved sign
is fair after conditioning. The actual limiting probability of
either sign at each such input is consequently at most
$1/2+\eta<2/3$. Every observed input has correct-label probability
one. Conditional convergence and the strict separation from $2/3$
then give the first assertion. Averaging over the finite set of
sample sequences is legitimate; each input is unobserved with
probability $(1-1/M)^n$. Solving the resulting equality proves the
sample bound.
\end{proof}

This theorem lets $d$ and the chosen depth vary after the width
limit. It does not prove an exponential bound for depth fixed
independently of $d$, nor for $0<\gamma<1$, where the effective
Gaussian margin diverges and a fixed near-independence estimate
does not control the zero-temperature posterior.

The same confidence conclusion sharpens the separate NNGP note
\cite{nngp-note}. That note uses first-layer variance $1/(d+1)$,
later hidden-layer and readout variances $1/(vN)$, where
$v=\mathbb E\tanh^2(G)$, and depth $O(d)$.
Its uniform Gaussian comparison gives
$|p(x\mid S,f)-1/2|\le\eta$ at every unseen input.
Taking $\eta<1/6$ makes confidence failure exactly the unseen
fraction, not merely close to it. Its expected value is again
\eqref{eq:deep-exact-confidence}. That is a concrete exponential
result under the requested confidence criterion, with the specified
layer gains and depth; it does not extend to every fixed-depth
kernel by changing the prediction rule.


\subsection{A small parity eigenvalue does not by itself force slow learning}

\begin{proposition}[A finite-rank Gaussian counterexample]
\label{prop:rank-two-counterexample}
Fix a nonconstant parity $f$ and $0<\lambda<1$. Consider the Gaussian
kernel
\[
 K_\lambda(x,y)=(1-\lambda)+\lambda f(x)f(y).
\]
It has normalized parity eigenvalue $\lambda$. Nevertheless, for its
Gaussian sign learner, either with any positive all-input margin or
without that margin, expected confidence-failure, posterior, and majority errors after
$n\ge1$ independent uniform labeled examples are at most $2^{-n}$.
The same bound holds for every zero-temperature limit obtained by
shrinking this Gaussian's overall amplitude while keeping the margin
fixed.
\end{proposition}

\begin{proof}
Write $G(x)=A+Bf(x)$ with independent
$A\sim\mathcal N(0,1-\lambda)$ and $B\sim\mathcal N(0,\lambda)$.
There are only two latent values, $A+B$ and $A-B$. With or without the
global margin event, the only possible sign functions are
$1,-1,f,-f$, each with positive probability. Once the training sample
contains both target labels, $f$ is the unique compatible sign
function. Otherwise the compatible signs are $f$ and the observed
constant, which disagree on half the domain. Every input on the
agreeing half has correct-label probability one; hence confidence
failure, posterior error, and majority error are all at most $1/2$.
The probability that all $n$ target
labels coincide is $2^{1-n}$. Multiplication yields the claimed bound,
uniformly in the Gaussian amplitude, so it persists in any of the
specified limits.
\end{proof}

This is a counterexample to an eigenvalue-only inference, not a claim
about the iterated-arcsine kernel~\eqref{eq:arcsine-kernel}. Its parity
eigenvalue can be arbitrarily small while learning remains fast. The
counterexample also shows why even an exact, very unfavorable parity
prior rate such as~\eqref{eq:parity-kernel-energy} does not alone yield
a lower bound on examples: a few observations can remove all
competitors with larger prior mass.

\section{The unresolved exponential parity question}

For the precise architecture~\eqref{eq:model-variances}, a fixed kernel
does govern every fixed $\gamma<1$. This follows from exponential
freezing, with no substitution of a regression likelihood for the
actual sign likelihood. Three statistical limits result:
\[
\begin{array}{c|l}
 0<\gamma<1 & \text{minimum kernel energy over compatible full-margin orthants},\\
 \gamma=0 & \text{Gaussian signs conditioned on every margin},\\
 \gamma<0 & \text{Gaussian signs conditioned on observed signs alone}.
\end{array}
\]
The freezing proof separates the scales $N$ and $N^\gamma$. At
$\gamma=1$ they coincide, so the argument cannot discard feature
deviations at the scale relevant to learning.

Note added (October 2026): the later paper \emph{Exponential parity sample complexity
below the mean-field scaling} proves an $\exp(\Omega(k))$ confidence lower bound
for full $k$-bit parity ($d=k$) at every fixed $\gamma<1$ and fixed depth, in the
width-first limit. For sparse parity with $0<\gamma<1$, the paper \emph{Exponential sample
lower bounds for sparse parity under a globally separated neural prior} proves an
$\exp(\Omega(k))$ bound; the natural $\exp\{\Omega(k\log(ed/k))\}$ scale at fixed depth
remains open. The following paragraph describes
the state of the argument when this paper was written.

The statements proved here do not establish a uniform
$\exp(\Omega(k))$ sample lower bound for $k$-bit parity throughout
$\gamma<1$. For positive $\gamma$ the missing step is a quantitative
comparison of the minimum energies of compatible sign patterns.
For zero or negative $\gamma$ at fixed depth, it is a quantitative
comparison of the relevant Gaussian orthant probabilities; the
depth-growing subfamily in Theorem~\ref{thm:deep-confidence} is
settled here. A regression mean's small
correlation with parity, or a small parity kernel eigenvalue, does
not supply either comparison. A majority-vote lower bound requires
additional control beyond posterior error. These are mathematical
gaps in an attempted deduction, rather than a proved counterexample
for the specified neural-network kernel.

\begin{thebibliography}{9}
\bibitem{notes}
D. Vaintrob and K. H\"anni.
\emph{A Generalization Result for Overparametrized Neural Net Bayesian Learning}.
Lecture notes, 2026.
\url{https://kaarelh.github.io/ai/nnbayes-generalization.pdf}.
\bibitem{nngp-note}
D. Vaintrob, K. H\"anni, and ChatGPT.
\emph{Small circuits under Gaussian-process classification:
what sample complexity follows from the actual posterior?}
Unpublished companion note, October 15, 2026.
\end{thebibliography}
\end{document}
